\documentclass[12pt,a4paper]{article}
\usepackage[T2A]{fontenc}
\usepackage[utf8]{inputenc}
\usepackage[english,russian]{babel}
\usepackage{amsmath,amssymb,amsfonts,mathtools}
\usepackage{geometry}
\usepackage{xcolor}
\usepackage{graphicx}
\usepackage{tikz}
\usepackage{enumitem}
\usepackage{multicol}
\usepackage{array,tabularx,booktabs}
\usepackage{fancyhdr}
\usepackage{lastpage}
\usepackage{microtype}
\usepackage{hyperref}

\geometry{
  left=20mm,
  right=20mm,
  top=20mm,
  bottom=22mm
}

\definecolor{accent}{HTML}{49677D}
\definecolor{accentlight}{HTML}{EAF0F4}
\definecolor{graphite}{HTML}{2F3438}
\definecolor{muted}{HTML}{6F7A82}

\hypersetup{
  colorlinks=true,
  linkcolor=accent,
  urlcolor=accent,
  pdftitle={Минимально необходимые темы математики 5 класса для освоения курса 6 класса},
  pdfauthor={Лёвин Артём Александрович}
}

\pagestyle{fancy}
\fancyhf{}
\fancyhead[L]{\small\textcolor{muted}{Математика: мост из 5 в 6 класс}}
\fancyhead[R]{\small\textcolor{muted}{Ксения Васильченко}}
\fancyfoot[C]{\small\textcolor{muted}{\thepage\ из \pageref{LastPage}}}
\renewcommand{\headrulewidth}{0.3pt}
\renewcommand{\footrulewidth}{0pt}
\setlength{\headheight}{15pt}

\setlength{\parindent}{0pt}
\setlength{\parskip}{0.55em}
\setlist[itemize]{leftmargin=1.4em,itemsep=0.25em,topsep=0.25em}
\setlist[enumerate]{leftmargin=1.8em,itemsep=0.45em,topsep=0.35em}

\newcommand{\topicline}{\par\vspace{0.25em}\noindent\textcolor{accent}{\rule{\textwidth}{0.5pt}}\par\vspace{0.35em}}
\newcommand{\term}[1]{\textbf{\textcolor{graphite}{#1}}}
\newcommand{\algorithm}[1]{%
  \par\smallskip
  {\color{accent}\textbf{Алгоритм.}} #1\par\smallskip
}
\newcommand{\attention}[1]{%
  \par\smallskip
  {\color{accent}\textbf{Обратите внимание.}} #1\par\smallskip
}
\newcommand{\mistake}[1]{%
  \par\smallskip
  {\color{accent}\textbf{Типичная ошибка.}} #1\par\smallskip
}
\newcommand{\selfcheck}[1]{%
  \par\smallskip
  {\color{accent}\textbf{Проверка себя.}} #1\par\smallskip
}
\newcommand{\memo}[1]{%
  \par\smallskip
  \begingroup
  \setlength{\fboxsep}{7pt}%
  \colorbox{accentlight}{\parbox{0.94\textwidth}{\textbf{Краткая памятка.} #1}}%
  \endgroup
  \par\smallskip
}
\newcommand{\trainingtitle}{\subsection*{Тренировка: 10 заданий}}
\newcommand{\answerline}[1]{\textbf{#1}}

\begin{document}

\begin{titlepage}
\thispagestyle{empty}
\begin{tikzpicture}[remember picture,overlay]
  \draw[accent!65,line width=0.8pt]
    ([xshift=-2mm,yshift=-31mm]current page.north west)
      .. controls +(45mm,-9mm) and +(-35mm,12mm) ..
    ([xshift=-12mm,yshift=-47mm]current page.north east);
\end{tikzpicture}

\vspace*{24mm}
{\large\bfseries\textcolor{accent}{ПОСОБИЕ ПО МАТЕМАТИКЕ}}\par
\vspace{11mm}
{\Huge\bfseries\textcolor{graphite}{Минимально необходимые темы\\[2mm]математики 5 класса}}\par
\vspace{3mm}
{\LARGE\textcolor{accent}{для уверенного освоения курса 6 класса}}\par
\vspace{12mm}
{\large Уровень: 5 класс, изучение новых тем}\par
\vspace{4mm}
{\large Дата: 04.09.26}\par

\vfill
{\small\textcolor{muted}{Лёвин Артём Александрович эксклюзивно для Ксения Васильченко}}\par
\end{titlepage}

\tableofcontents
\newpage

\section*{Как устроено пособие}
\addcontentsline{toc}{section}{Как устроено пособие}

Это пособие не пытается повторить весь курс 5 класса. Его задача --- дать \textbf{минимальный, но достаточный фундамент}, на который в 6 классе будут опираться темы натуральных чисел, дробей, отношений, процентов, отрицательных чисел, буквенных выражений, задач и наглядной геометрии.

Материал выстроен по траектории
\[
\text{узнаю}\;\to\;\text{воспроизвожу}\;\to\;\text{выбираю метод}\;\to\;\text{комбинирую}\;\to\;\text{применяю}.
\]

После каждого смыслового блока дано ровно 10 заданий. Сначала идут прямые применения, затем комбинированные и перенос метода в новую ситуацию. Ответы находятся в конце пособия.

\section{Карта темы и предварительные знания}

\subsection*{Что желательно уметь до начала}
Из курса начальной школы достаточно уверенно владеть:
\begin{itemize}
  \item таблицей умножения;
  \item письменным сложением и вычитанием;
  \item смыслом умножения и деления;
  \item простейшими единицами длины, массы и времени;
  \item понятиями прямоугольника и квадрата.
\end{itemize}

Если какой-то навык пока нестабилен, это не мешает начать: в пособии ключевые алгоритмы повторяются в нужном месте.

\subsection*{Восемь опорных блоков}
\begin{enumerate}
  \item Натуральные числа: запись, сравнение, числовая прямая, округление.
  \item Арифметические действия, свойства, выражения, степень, деление с остатком.
  \item Делители, кратные, простые числа и признаки делимости.
  \item Обыкновенные дроби: смысл, смешанные числа, основное свойство, сравнение.
  \item Действия с обыкновенными дробями и основные задачи на дроби.
  \item Десятичные дроби: запись, сравнение, округление и вычисления.
  \item Неизвестный компонент, формулы величин и текстовые задачи.
  \item Наглядная геометрия: углы, периметр, площадь и объём.
\end{enumerate}

\subsection*{Что сознательно не входит в этот минимум}
\begin{itemize}
  \item римская нумерация и подробное изучение непозиционных систем счисления;
  \item сложные переборные и комбинаторные задачи;
  \item подробная работа с развёртками многогранников и специальными построениями;
  \item \textbf{НОД и НОК} --- это уже систематическая тема 6 класса;
  \item \textbf{отношения, пропорции, масштаб и проценты} --- они изучаются как новые темы 6 класса;
  \item \textbf{положительные и отрицательные числа, модуль, координатная плоскость} --- новые темы 6 класса.
\end{itemize}

\attention{Основы делимости всё же включены: делители, кратные, простые числа и признаки делимости нужны как стартовая площадка для НОД, НОК и дробей в 6 классе.}

\newpage

\section{Натуральные числа: запись, сравнение и округление}

\subsection*{1. Что такое натуральные числа}
Натуральные числа используются при счёте:
\[
1,2,3,4,\ldots
\]
Число $0$ обозначает отсутствие объектов и рассматривается вместе с натуральными числами в арифметике 5 класса.

\term{Десятичная позиционная запись} означает, что значение цифры зависит от её места. Например,
\[
508\,307=5\cdot100000+0\cdot10000+8\cdot1000+3\cdot100+0\cdot10+7.
\]
То есть
\[
508\,307=500\,000+8\,000+300+7.
\]

\subsection*{2. Сравнение натуральных чисел}
Если числа имеют разное количество цифр, больше то, в котором цифр больше. Если количество цифр одинаково, сравниваем слева направо до первой различающейся цифры.

\textbf{Пример.} Сравним $405\,090$ и $405\,900$. Первые три цифры совпадают. Затем сравниваем сотни: $0<9$, поэтому
\[
405\,090<405\,900.
\]

\subsection*{3. Числовая прямая}
На числовой прямой выбирают начало отсчёта $O$, направление и единичный отрезок. Чем правее расположена точка, тем больше соответствующее число.

\begin{center}
\begin{tikzpicture}[scale=0.9]
  \draw[->,line width=0.7pt] (0,0)--(9.4,0);
  \foreach \x in {0,...,8}{
    \draw (\x,-0.10)--(\x,0.10);
    \node[below] at (\x,-0.12) {\small \x};
  }
  \fill[accent] (6,0) circle (2pt);
  \node[above,accent] at (6,0.12) {$A$};
\end{tikzpicture}
\end{center}

Здесь точке $A$ соответствует число $6$.

\subsection*{4. Округление}
Округление заменяет число близким числом с нулями в младших разрядах.

\algorithm{Чтобы округлить до заданного разряда: 1) найдите цифру этого разряда; 2) посмотрите на цифру справа; 3) если она $0,1,2,3,4$, выбранную цифру не меняем; если $5,6,7,8,9$, увеличиваем на $1$; 4) все цифры справа заменяем нулями.}

\textbf{Пример.} Округлим $87\,649$ до сотен. Цифра сотен --- $6$, справа стоит $4$, поэтому
\[
87\,649\approx87\,600.
\]
До тысяч:
\[
87\,649\approx88\,000,
\]
потому что справа от цифры тысяч $7$ стоит $6\ge5$.

\mistake{Нельзя смотреть на все цифры справа сразу. Решение об увеличении разряда принимает только ближайшая цифра справа.}

\selfcheck{После округления число должно быть близко к исходному. Например, $87\,649$ не может округлиться до $80\,000$ при округлении до тысяч: $88\,000$ гораздо ближе.}

\memo{Разрядная запись помогает понимать число; сравнение выполняем слева направо; на числовой прямой большее число правее; при округлении смотрим на одну цифру справа от нужного разряда.}

\trainingtitle
\begin{enumerate}
  \item Запишите цифрами число: семь миллионов сорок две тысячи пятнадцать.
  \item Представьте число $508\,307$ в виде суммы разрядных слагаемых.
  \item Сравните: $405\,090$ и $405\,900$.
  \item Расположите по возрастанию: $12\,050$, $1\,250$, $10\,250$, $120\,500$.
  \item Округлите $87\,649$ 1) до сотен; 2) до тысяч.
  \item Единичный отрезок на числовой прямой равен $2$ см. На каком расстоянии от нуля находится точка с координатой $7$?
  \item Какое натуральное число расположено строго между $9\,999$ и $10\,001$?
  \item Перечислите все натуральные $x$, для которых $3<x<8$.
  \item Выполните прикидку суммы $3\,987+6\,124$, округлив каждое слагаемое до тысяч.
  \item Число при округлении до сотен стало равно $12\,400$. Укажите наименьшее и наибольшее натуральные числа, которые могли быть исходными.
\end{enumerate}

\newpage

\section{Действия с натуральными числами и числовые выражения}

\subsection*{1. Четыре арифметических действия}
\begin{center}
\begin{tabularx}{\textwidth}{@{}>{\bfseries}l X X@{}}
\toprule
Действие & Компоненты & Связь для проверки \\
\midrule
Сложение & слагаемое $+$ слагаемое $=$ сумма & сумма $-$ слагаемое $=$ другое слагаемое \\
Вычитание & уменьшаемое $-$ вычитаемое $=$ разность & разность $+$ вычитаемое $=$ уменьшаемое \\
Умножение & множитель $\cdot$ множитель $=$ произведение & произведение $:$ множитель $=$ другой множитель \\
Деление & делимое $:$ делитель $=$ частное & частное $\cdot$ делитель $=$ делимое \\
\bottomrule
\end{tabularx}
\end{center}

\subsection*{2. Свойства, которые экономят вычисления}
Для сложения и умножения работают переместительное и сочетательное свойства:
\[
a+b=b+a,\qquad ab=ba,
\]
\[
(a+b)+c=a+(b+c),\qquad (ab)c=a(bc).
\]
Распределительное свойство:
\[
a(b+c)=ab+ac.
\]

\textbf{Пример.}
\begin{align*}
25\cdot37\cdot4&=(25\cdot4)\cdot37\\
&=100\cdot37\\
&=3700.
\end{align*}

Ещё один полезный приём:
\begin{align*}
49\cdot18+51\cdot18&=(49+51)\cdot18\\
&=100\cdot18=1800.
\end{align*}

\subsection*{3. Порядок действий}
\algorithm{Сначала выполняем действия в скобках. Затем степени. Потом умножение и деление слева направо. В конце сложение и вычитание слева направо.}

\textbf{Пример.}
\begin{align*}
720-6(35+15)&=720-6\cdot50\\
&=720-300\\
&=420.
\end{align*}

\subsection*{4. Степень}
Запись $a^n$ означает произведение $n$ одинаковых множителей $a$:
\[
a^n=\underbrace{a\cdot a\cdot\ldots\cdot a}_{n\text{ множителей}}.
\]
Например,
\[
3^4=3\cdot3\cdot3\cdot3=81.
\]

\subsection*{5. Деление с остатком}
Если число не делится нацело, используют запись
\[
a=bq+r,
\]
где $a$ --- делимое, $b$ --- делитель, $q$ --- неполное частное, $r$ --- остаток, причём
\[
0\le r<b.
\]

\textbf{Пример.} $785:23=34$ (ост. $3$), потому что
\[
785=23\cdot34+3.
\]

\mistake{Остаток не может быть равен делителю или быть больше него. Если это произошло, деление ещё не закончено.}

\selfcheck{Проверяйте сложные вычисления обратным действием и прикидкой. Если $9360:45$, результат должен быть около $9000:45=200$, поэтому ответ порядка $200$ реалистичен.}

\memo{Используйте свойства действий для удобных группировок; соблюдайте порядок действий; степень --- короткая запись повторного умножения; при делении с остатком всегда проверяйте $a=bq+r$ и $r<b$.}

\trainingtitle
\begin{enumerate}
  \item Вычислите: $3478+6522$.
  \item Вычислите: $9000-4685$.
  \item Вычислите: $125\cdot48$.
  \item Вычислите: $9360:45$.
  \item Найдите значение: $720-6(35+15)$.
  \item Вычислите удобным способом: $25\cdot37\cdot4$.
  \item Вычислите удобным способом: $49\cdot18+51\cdot18$.
  \item Найдите значение: $3^4+2^5$.
  \item Выполните деление с остатком: $785:23$.
  \item Найдите $x$: $(x+135)\cdot4=1000$.
\end{enumerate}

\newpage

\section{Делители, кратные и признаки делимости}

\subsection*{1. Делитель и кратное}
Если натуральное число $a$ делится на $b$ без остатка, то $b$ называется \term{делителем} числа $a$, а $a$ --- \term{кратным} числа $b$.

Например, делители числа $18$:
\[
1,2,3,6,9,18.
\]
Числа $18,36,54,72,\ldots$ --- кратные $18$.

\subsection*{2. Простые и составные числа}
\term{Простое число} имеет ровно два натуральных делителя: $1$ и само число. Например, $2,3,5,7,11,13,17$.

\term{Составное число} имеет больше двух делителей. Например, $12$ делится на $1,2,3,4,6,12$.

Число $1$ не является ни простым, ни составным.

\subsection*{3. Разложение на простые множители}
Число можно последовательно делить на простые делители.

\textbf{Пример.}
\begin{align*}
84&=2\cdot42\\
  &=2\cdot2\cdot21\\
  &=2^2\cdot3\cdot7.
\end{align*}

\subsection*{4. Признаки делимости}
\begin{center}
\begin{tabularx}{\textwidth}{@{}c X@{}}
\toprule
На число & Признак \\
\midrule
$2$ & последняя цифра чётная: $0,2,4,6,8$ \\
$5$ & последняя цифра $0$ или $5$ \\
$10$ & последняя цифра $0$ \\
$3$ & сумма цифр делится на $3$ \\
$9$ & сумма цифр делится на $9$ \\
\bottomrule
\end{tabularx}
\end{center}

\textbf{Пример.} Для числа $7\,326$ сумма цифр равна
\[
7+3+2+6=18.
\]
Значит, $7\,326$ делится и на $3$, и на $9$. Последняя цифра $6$, поэтому число также делится на $2$.

\attention{В 6 классе на этой основе появятся НОД и НОК. Здесь достаточно уверенно видеть делители, кратные и применять признаки делимости.}

\mistake{Признак делимости на $3$ относится к сумме цифр, а не к последней цифре. Например, $123$ делится на $3$, потому что $1+2+3=6$.}

\memo{Делитель делит число нацело; кратные получаются умножением; простое число имеет ровно два делителя; для $3$ и $9$ проверяем сумму цифр, для $2,5,10$ --- последнюю цифру.}

\trainingtitle
\begin{enumerate}
  \item Выпишите все натуральные делители числа $18$.
  \item Запишите первые пять положительных кратных числа $7$.
  \item Из чисел $120$, $135$, $248$, $305$, $410$ выпишите 1) делящиеся на $2$; 2) на $5$; 3) на $10$.
  \item Из чисел $246$, $451$, $729$, $1008$ выпишите 1) делящиеся на $3$; 2) делящиеся на $9$.
  \item Разделите числа $2$, $17$, $21$, $29$, $35$, $51$ на простые и составные.
  \item Разложите $84$ на простые множители.
  \item Какую наименьшую цифру можно поставить вместо $*$ в числе $4*2$, чтобы число делилось на $3$?
  \item Найдите наибольшее трёхзначное число, которое делится и на $5$, и на $9$.
  \item Выпишите все числа от $1$ до $50$, которые делятся на $6$, но не делятся на $4$.
  \item Найдите наименьшее натуральное число, большее $100$, которое делится одновременно на $2$, $3$ и $5$.
\end{enumerate}

\newpage

\section{Обыкновенные дроби: смысл, преобразования и сравнение}

\subsection*{1. Смысл дроби}
Дробь
\[
\frac{a}{b},\qquad b\ne0,
\]
показывает, что целое разделили на $b$ равных частей и взяли $a$ таких частей. Число $a$ --- \term{числитель}, $b$ --- \term{знаменатель}.

Например, $\frac{3}{7}$ означает три части из семи равных частей.

\subsection*{2. Правильные, неправильные и смешанные дроби}
Если числитель меньше знаменателя, дробь правильная:
\[
\frac35<1.
\]
Если числитель не меньше знаменателя, дробь неправильная:
\[
\frac94>1,\qquad \frac77=1.
\]

Смешанное число состоит из целой и дробной частей:
\[
2\frac35=2+\frac35.
\]

\algorithm{Чтобы перевести $m\frac ab$ в неправильную дробь, умножьте целую часть $m$ на знаменатель $b$, прибавьте числитель $a$ и запишите результат над тем же знаменателем: $m\frac ab=\frac{mb+a}{b}$.}

\textbf{Пример.}
\[
2\frac35=\frac{2\cdot5+3}{5}=\frac{13}{5}.
\]

\algorithm{Чтобы выделить целую часть из неправильной дроби, разделите числитель на знаменатель с остатком. Частное станет целой частью, остаток --- числителем дробной части.}

\textbf{Пример.}
\[
\frac{17}{4}=4\frac14,
\]
потому что $17=4\cdot4+1$.

\subsection*{3. Основное свойство дроби}
Если числитель и знаменатель умножить или разделить на одно и то же ненулевое число, значение дроби не изменится:
\[
\frac ab=\frac{ak}{bk}.
\]
Например,
\[
\frac34=\frac{18}{24}.
\]

\term{Сократить дробь} --- разделить числитель и знаменатель на общий делитель больше $1$.

\textbf{Пример.}
\[
\frac{18}{30}=\frac{18:6}{30:6}=\frac35.
\]

\subsection*{4. Сравнение дробей}
Если знаменатели одинаковые, больше дробь с большим числителем. Если числители одинаковые и дроби положительные, больше дробь с меньшим знаменателем.

В общем случае приводим дроби к общему знаменателю.

\textbf{Пример.}
\[
\frac58\quad\text{и}\quad\frac34=\frac68,
\]
поэтому
\[
\frac58<\frac34.
\]

\mistake{Нельзя складывать или сравнивать дроби, механически складывая знаменатели. Знаменатель показывает размер частей; сначала части должны стать одинакового размера.}

\memo{Смешанное число переводим в неправильную дробь по формуле $\frac{mb+a}{b}$; основное свойство позволяет сокращать и приводить к новому знаменателю; для сравнения удобно получить общий знаменатель.}

\trainingtitle
\begin{enumerate}
  \item В дроби $\frac37$ назовите числитель и знаменатель и объясните их смысл.
  \item Какие из дробей $\frac35$, $\frac77$, $\frac94$, $\frac{11}{20}$ правильные, а какие неправильные?
  \item Переведите $2\frac35$ в неправильную дробь.
  \item Выделите целую часть: $\frac{17}{4}$.
  \item Запишите дробь, равную $\frac34$, со знаменателем $24$.
  \item Сократите дробь $\frac{18}{30}$.
  \item Сравните $\frac58$ и $\frac34$.
  \item Расположите по возрастанию: $\frac12$, $\frac23$, $\frac34$.
  \item Какая дробь находится ровно посередине между $\frac14$ и $\frac34$ на числовой прямой?
  \item Выпишите все дроби со знаменателем $12$, которые строго больше $\frac12$, но строго меньше $\frac34$.
\end{enumerate}

\newpage

\section{Действия с обыкновенными дробями и задачи на дроби}

\subsection*{1. Сложение и вычитание}
При одинаковых знаменателях складываем или вычитаем числители, знаменатель сохраняем:
\[
\frac ab+\frac cb=\frac{a+c}{b}.
\]

При разных знаменателях сначала приводим дроби к общему знаменателю.

\textbf{Пример.}
\begin{align*}
\frac23+\frac14&=\frac8{12}+\frac3{12}\\
&=\frac{11}{12}.
\end{align*}

\subsection*{2. Умножение дробей}
\[
\frac ab\cdot\frac cd=\frac{ac}{bd}.
\]
До умножения полезно выполнить сокращение крест-накрест.

\textbf{Пример.}
\begin{align*}
\frac35\cdot\frac{10}{9}
&=\frac{3\cdot10}{5\cdot9}\\
&=\frac23.
\end{align*}

\subsection*{3. Деление дробей}
Две дроби \term{взаимно обратны}, если их произведение равно $1$. Для $\frac ab$ обратная дробь --- $\frac ba$.

\algorithm{Чтобы разделить на дробь, умножьте на обратную дробь: $\frac ab:\frac cd=\frac ab\cdot\frac dc$.}

\textbf{Пример.}
\begin{align*}
\frac7{12}:\frac{14}{15}
&=\frac7{12}\cdot\frac{15}{14}\\
&=\frac58.
\end{align*}

\subsection*{4. Часть от целого и целое по его части}
Чтобы найти дробь от числа, число умножают на дробь.

\textbf{Пример.} Найти $\frac38$ от $64$:
\[
64\cdot\frac38=24.
\]

Чтобы найти целое по известной части, значение части делят на соответствующую дробь.

\textbf{Пример.} Если $\frac57$ числа равны $45$, то
\[
45:\frac57=45\cdot\frac75=63.
\]

\mistake{При делении дробей переворачивается только \emph{делитель}, а не обе дроби.}

\selfcheck{Если ищем часть от положительного числа и дробь меньше $1$, ответ должен быть меньше целого. Если $\frac38$ от $64$ получилось больше $64$, в вычислении есть ошибка.}

\memo{Сложение и вычитание требуют общего знаменателя; при умножении перемножаем числители и знаменатели; деление заменяем умножением на обратную дробь; часть от целого --- умножение, целое по части --- деление на дробь.}

\trainingtitle
\begin{enumerate}
  \item Вычислите: $\frac38+\frac28$.
  \item Вычислите: $\frac79-\frac49$.
  \item Вычислите: $\frac23+\frac14$.
  \item Вычислите: $\frac56-\frac14$.
  \item Вычислите: $\frac35\cdot\frac{10}{9}$.
  \item Вычислите: $\frac7{12}:\frac{14}{15}$.
  \item Вычислите: $2\frac13+1\frac56$.
  \item Найдите $\frac38$ от $64$.
  \item Найдите число, если $\frac57$ этого числа равны $45$.
  \item Вычислите: $\left(\frac34-\frac16\right):\frac78$.
\end{enumerate}

\newpage

\section{Десятичные дроби}

\subsection*{1. Что означает десятичная запись}
Десятичная дробь использует разряды десятых, сотых, тысячных и так далее.

Например,
\[
5{,}037=5+\frac{0}{10}+\frac3{100}+\frac7{1000}.
\]
Это число читается: «пять целых тридцать семь тысячных».

Десятичную дробь можно представить обыкновенной:
\[
0{,}125=\frac{125}{1000}=\frac18.
\]

\subsection*{2. Сравнение}
Сначала сравниваем целые части. Если они равны, сравниваем цифры после запятой по разрядам. Можно дописывать справа нули:
\[
4{,}08=4{,}080<4{,}800=4{,}8.
\]

\subsection*{3. Округление}
Правило такое же, как для натуральных чисел. Например,
\[
18{,}376\approx18{,}4\quad\text{до десятых},
\]
\[
18{,}376\approx18{,}38\quad\text{до сотых}.
\]

\subsection*{4. Сложение и вычитание}
\algorithm{Записывайте числа так, чтобы запятая была под запятой. При необходимости дописывайте нули справа. Затем выполняйте обычное сложение или вычитание по разрядам.}

\textbf{Пример.}
\[
12{,}750+8{,}906=21{,}656.
\]

\subsection*{5. Умножение}
Сначала умножаем числа, временно не учитывая запятые. Затем отделяем справа столько знаков, сколько их было после запятых в обоих множителях вместе.

\textbf{Пример.}
\[
3{,}6\cdot2{,}5=9{,}00=9.
\]

\subsection*{6. Деление}
Если делитель десятичный, переносим запятую в делимом и делителе на одинаковое число разрядов так, чтобы делитель стал натуральным.

\textbf{Пример.}
\[
15{,}75:0{,}5=157{,}5:5=31{,}5.
\]

\mistake{При сложении и вычитании выравнивают \emph{запятые}, а не последние цифры. При умножении, наоборот, запятые сначала не выравнивают.}

\selfcheck{Оцените порядок результата. $3{,}6\cdot2{,}5$ должно быть около $4\cdot2{,}5=10$, значит ответ $9$ реалистичен, а $90$ --- нет.}

\memo{Сравниваем разряды; округляем по соседней цифре справа; при сложении и вычитании ставим запятые друг под другом; при делении на десятичную дробь сначала превращаем делитель в натуральное число.}

\trainingtitle
\begin{enumerate}
  \item Запишите десятичной дробью: пять целых тридцать семь тысячных.
  \item Представьте $0{,}125$ в виде несократимой обыкновенной дроби.
  \item Сравните: $4{,}08$ и $4{,}8$.
  \item Расположите по возрастанию: $0{,}9$, $0{,}09$, $0{,}909$, $0{,}099$.
  \item Округлите $18{,}376$ 1) до десятых; 2) до сотых.
  \item Вычислите: $12{,}75+8{,}906$.
  \item Вычислите: $40-7{,}385$.
  \item Вычислите: $3{,}6\cdot2{,}5$.
  \item Вычислите: $15{,}75:0{,}5$.
  \item Вычислите: $2{,}4(3{,}5-1{,}25)+0{,}6$.
\end{enumerate}

\newpage

\section{Неизвестный компонент, величины и текстовые задачи}

\subsection*{1. Буква как неизвестное число}
Буква может обозначать неизвестное число. Для простых уравнений удобно не «переносить через знак равно», а использовать связь между компонентами действий.

\begin{center}
\begin{tabularx}{\textwidth}{@{}X X@{}}
\toprule
Если известно & Как найти неизвестное \\
\midrule
$x+a=b$ & $x=b-a$ \\
$x-a=b$ & $x=b+a$ \\
$a-x=b$ & $x=a-b$ \\
$ax=b$ & $x=b:a$ \\
$x:a=b$ & $x=ab$ \\
$a:x=b$ & $x=a:b$ \\
\bottomrule
\end{tabularx}
\end{center}

\textbf{Пример.}
\begin{align*}
(x+15)\cdot7&=280,\\
x+15&=280:7=40,\\
x&=40-15=25.
\end{align*}
Проверка:
\[
(25+15)\cdot7=40\cdot7=280.
\]

\subsection*{2. Три главные зависимости}
Для движения:
\[
S=vt,
\]
где $S$ --- расстояние, $v$ --- скорость, $t$ --- время. Отсюда
\[
v=\frac St,\qquad t=\frac Sv.
\]

Для покупок:
\[
C=pn,
\]
где $C$ --- стоимость, $p$ --- цена одной единицы, $n$ --- количество.

\subsection*{3. Единицы измерения}
Перед вычислениями величины должны быть выражены в согласованных единицах.

Полезные связи:
\[
1\text{ км}=1000\text{ м},\qquad 1\text{ м}=100\text{ см},
\]
\[
1\text{ кг}=1000\text{ г},\qquad 1\text{ ч}=60\text{ мин}.
\]

\textbf{Пример.} $2$ ч $15$ мин:
\[
2\cdot60+15=135\text{ мин}.
\]

\subsection*{4. Арифметический способ решения задач}
До систематического изучения алгебры многие задачи удобно решать рассуждением по действиям.

\textbf{Пример.} На двух полках $86$ книг. На второй на $18$ книг больше, чем на первой. Если убрать «лишние» $18$ книг со второй полки, останутся две равные части:
\[
86-18=68,
\]
\[
68:2=34.
\]
На первой полке $34$ книги, на второй
\[
34+18=52.
\]

\algorithm{Для текстовой задачи: 1) выпишите данные и единицы; 2) определите, что неизвестно; 3) найдите зависимость между величинами; 4) составьте последовательность действий; 5) вычислите; 6) проверьте смысл ответа и единицы.}

\mistake{Нельзя подставлять в одну формулу несогласованные единицы. Например, скорость в км/ч и время в минутах сначала требуют перевода времени в часы или скорости в км/мин.}

\memo{Уравнения решаем через обратные действия и проверяем подстановкой; движение: $S=vt$; покупки: стоимость = цена $\times$ количество; в задачах сначала приводим единицы к совместимому виду.}

\trainingtitle
\begin{enumerate}
  \item Решите уравнение: $x+378=950$.
  \item Решите уравнение: $1200-x=485$.
  \item Решите уравнение: $36x=900$.
  \item Решите уравнение: $x:24=15$.
  \item Решите уравнение: $(x+15)\cdot7=280$.
  \item Автомобиль ехал со скоростью $72$ км/ч в течение $3{,}5$ ч. Какое расстояние он проехал?
  \item Поезд прошёл $240$ км за $4$ ч с постоянной скоростью. Найдите скорость.
  \item $8$ одинаковых тетрадей стоят $56$ рублей. Сколько стоят $13$ таких тетрадей?
  \item Переведите: 1) $3{,}6$ км в метры; 2) $2$ ч $15$ мин в минуты; 3) $4500$ г в килограммы.
  \item На двух полках $86$ книг. На второй на $18$ книг больше, чем на первой. Сколько книг на каждой полке?
\end{enumerate}

\newpage

\section{Наглядная геометрия: длина, угол, площадь и объём}

\subsection*{1. Базовые фигуры}
\term{Прямая} не имеет ни начала, ни конца. \term{Луч} имеет начало и не имеет конца в одном направлении. \term{Отрезок} ограничен двумя концами.

Если точка $C$ лежит между $A$ и $B$, то
\[
AB=AC+CB.
\]

\subsection*{2. Угол}
Угол образован двумя лучами с общим началом. Величину угла измеряют в градусах.

\begin{center}
\begin{tabular}{@{}ll@{}}
\toprule
Вид угла & Величина \\
\midrule
острый & меньше $90^\circ$ \\
прямой & $90^\circ$ \\
тупой & больше $90^\circ$, но меньше $180^\circ$ \\
развёрнутый & $180^\circ$ \\
\bottomrule
\end{tabular}
\end{center}

\subsection*{3. Периметр и площадь}
Периметр --- сумма длин всех сторон.

Для прямоугольника со сторонами $a$ и $b$:
\[
P=2(a+b),\qquad S=ab.
\]
Для квадрата со стороной $a$:
\[
P=4a,\qquad S=a^2.
\]

\textbf{Пример.} Прямоугольник $8$ см на $5$ см:
\[
P=2(8+5)=26\text{ см},
\]
\[
S=8\cdot5=40\text{ см}^2.
\]

\subsection*{4. Составные фигуры}
Если фигура составлена из прямоугольников, её площадь можно найти сложением площадей частей или вычитанием вырезанной части из большой.

\textbf{Пример.} Из прямоугольника $10\times8$ вырезали прямоугольник $3\times2$:
\[
S=10\cdot8-3\cdot2=80-6=74.
\]

\subsection*{5. Окружность и круг}
\term{Радиус} соединяет центр окружности с её точкой. \term{Диаметр} проходит через центр и соединяет две точки окружности:
\[
d=2r.
\]

\subsection*{6. Объём прямоугольного параллелепипеда}
Для рёбер $a,b,c$:
\[
V=abc.
\]
Для куба:
\[
V=a^3.
\]

Полезно помнить:
\[
1\text{ дм}^3=1000\text{ см}^3.
\]

\mistake{Единицы площади и объёма нельзя переводить как единицы длины. Например, $1\text{ м}^2=10\,000\text{ см}^2$, потому что $100\cdot100=10\,000$.}

\selfcheck{Периметр измеряется в единицах длины, площадь --- в квадратных единицах, объём --- в кубических. Единица в ответе помогает обнаружить ошибку формулы.}

\memo{Отрезки складываются по длине; углы различаются относительно $90^\circ$ и $180^\circ$; $P_{\text{прям.}}=2(a+b)$, $S_{\text{прям.}}=ab$, $V_{\text{паралл.}}=abc$, диаметр равен двум радиусам.}

\trainingtitle
\begin{enumerate}
  \item Точка $C$ лежит между $A$ и $B$. $AB=7{,}5$ см, $AC=2{,}8$ см. Найдите $CB$.
  \item Определите вид каждого угла: $38^\circ$, $90^\circ$, $127^\circ$, $180^\circ$.
  \item Найдите периметр прямоугольника со сторонами $8$ см и $5$ см.
  \item Найдите площадь прямоугольника со сторонами $12$ см и $7$ см.
  \item Площадь квадрата равна $81$ см$^2$. Найдите сторону и периметр квадрата.
  \item Из прямоугольника $10$ см $\times$ $8$ см вырезали прямоугольник $3$ см $\times$ $2$ см. Найдите площадь оставшейся фигуры.
  \item Радиус окружности равен $4{,}5$ см. Найдите диаметр.
  \item Найдите объём прямоугольного параллелепипеда с измерениями $6$ см, $4$ см и $3$ см.
  \item Переведите: 1) $3$ м$^2$ в см$^2$; 2) $2{,}5$ дм$^3$ в см$^3$.
  \item Периметр прямоугольника равен $54$ см, его длина --- $16$ см. Найдите ширину и площадь.
\end{enumerate}

\newpage

\section{Связи между темами: как выбирать метод}

В 6 классе задачи часто требуют использование сочетания нескольких навыков. Поэтому полезно научиться распознавать «сигналы» метода.

\begin{center}
\begin{tabularx}{\textwidth}{@{}>{\bfseries}X X@{}}
\toprule
Что видно в условии & Что вспомнить \\
\midrule
Большие натуральные числа, скобки & порядок действий, свойства, прикидка \\
Нужно понять, делится ли число & признаки делимости \\
Обыкновенные дроби разных знаменателей & общий знаменатель \\
Деление на дробь & умножение на обратную \\
«Часть от числа» & умножить на дробь \\
«Известна часть, найти целое» & разделить на дробь \\
Десятичные дроби & положение запятой и оценка результата \\
Скорость, время, расстояние & $S=vt$ \\
Цена, количество, стоимость & $C=pn$ \\
Прямоугольник или параллелепипед & формулы $P$, $S$, $V$ \\
\bottomrule
\end{tabularx}
\end{center}

\subsection*{Проверка готовности к 6 классу}
Можно считать фундамент достаточным, если без подсказки удаётся:
\begin{itemize}
  \item вычислять выражения с натуральными числами и скобками;
  \item быстро применять признаки делимости на $2,3,5,9,10$;
  \item переводить смешанные числа в неправильные дроби и обратно;
  \item складывать, вычитать, умножать и делить простые дроби;
  \item уверенно выполнять четыре действия с десятичными дробями;
  \item решать простые задачи на часть и целое, движение и покупки;
  \item пользоваться формулами периметра, площади и объёма;
  \item выполнять прикидку результата и проверять единицы измерения.
\end{itemize}

\newpage

\section{Итоговая обзорная контрольная работа}

Контрольная охватывает весь необходимый фундамент. Рекомендуемый режим: без подсказок, с обязательной записью промежуточных действий.

\begin{enumerate}
  \item \textbf{Базовый уровень.} Округлите число $607\,948$ до тысяч.
  \item \textbf{Базовый уровень.} Вычислите:
  \[
  8400:35+27\cdot16.
  \]
  \item \textbf{Базовый уровень.} Найдите наименьшее натуральное число, большее $200$, которое делится одновременно на $2$, $3$ и $5$.
  \item \textbf{Средний уровень.} Сократите дробь $\frac{42}{56}$.
  \item \textbf{Средний уровень.} Вычислите:
  \[
  \frac56-\frac13+\frac14.
  \]
  \item \textbf{Средний уровень.} Вычислите:
  \[
  7{,}2:0{,}9+3{,}75.
  \]
  \item \textbf{Средний уровень.} Решите уравнение:
  \[
  (x-18)\cdot5=210.
  \]
  \item \textbf{Повышенный уровень.} Автомобиль ехал со скоростью $64$ км/ч в течение $2{,}5$ ч. Какое расстояние он прошёл?
  \item \textbf{Повышенный уровень.} Площадь прямоугольника равна $96$ см$^2$, ширина --- $8$ см. Найдите длину и периметр.
  \item \textbf{Комплексное задание.} В баке было $120$ л воды. Сначала использовали $\frac38$ всей воды, затем ещё $25{,}5$ л. Сколько литров воды осталось?
\end{enumerate}

\newpage

\section{Ответы}

\subsection*{К блоку 1}
\begin{enumerate}
  \item $7\,042\,015$.
  \item $500\,000+8\,000+300+7$.
  \item $405\,090<405\,900$.
  \item $1\,250<10\,250<12\,050<120\,500$.
  \item 1) $87\,600$; 2) $88\,000$.
  \item $14$ см.
  \item $10\,000$.
  \item $4,5,6,7$.
  \item Примерно $10\,000$.
  \item От $12\,350$ до $12\,449$; наименьшее $12\,350$, наибольшее $12\,449$.
\end{enumerate}

\subsection*{К блоку 2}
\begin{enumerate}
  \item $10\,000$.
  \item $4315$.
  \item $6000$.
  \item $208$.
  \item $420$.
  \item $3700$.
  \item $1800$.
  \item $113$.
  \item $34$ (ост. $3$).
  \item $x=115$.
\end{enumerate}

\subsection*{К блоку 3}
\begin{enumerate}
  \item $1,2,3,6,9,18$.
  \item $7,14,21,28,35$.
  \item 1) $120,248,410$; 2) $120,135,305,410$; 3) $120,410$.
  \item 1) $246,729,1008$; 2) $729,1008$.
  \item Простые: $2,17,29$; составные: $21,35,51$.
  \item $84=2^2\cdot3\cdot7$.
  \item $0$; получается $402$.
  \item $990$.
  \item $6,18,30,42$.
  \item $120$.
\end{enumerate}

\subsection*{К блоку 4}
\begin{enumerate}
  \item Числитель $3$, знаменатель $7$: взяты $3$ из $7$ равных частей.
  \item Правильные: $\frac35$, $\frac{11}{20}$; неправильные: $\frac77$, $\frac94$.
  \item $\frac{13}{5}$.
  \item $4\frac14$.
  \item $\frac{18}{24}$.
  \item $\frac35$.
  \item $\frac58<\frac34$.
  \item $\frac12<\frac23<\frac34$.
  \item $\frac12$.
  \item $\frac7{12}$ и $\frac8{12}$.
\end{enumerate}

\subsection*{К блоку 5}
\begin{enumerate}
  \item $\frac58$.
  \item $\frac13$.
  \item $\frac{11}{12}$.
  \item $\frac7{12}$.
  \item $\frac23$.
  \item $\frac58$.
  \item $4\frac16$.
  \item $24$.
  \item $63$.
  \item $\frac23$.
\end{enumerate}

\subsection*{К блоку 6}
\begin{enumerate}
  \item $5{,}037$.
  \item $\frac18$.
  \item $4{,}08<4{,}8$.
  \item $0{,}09<0{,}099<0{,}9<0{,}909$.
  \item 1) $18{,}4$; 2) $18{,}38$.
  \item $21{,}656$.
  \item $32{,}615$.
  \item $9$.
  \item $31{,}5$.
  \item $6$.
\end{enumerate}

\subsection*{К блоку 7}
\begin{enumerate}
  \item $x=572$.
  \item $x=715$.
  \item $x=25$.
  \item $x=360$.
  \item $x=25$.
  \item $252$ км.
  \item $60$ км/ч.
  \item $91$ рубль.
  \item 1) $3600$ м; 2) $135$ мин; 3) $4{,}5$ кг.
  \item На первой полке $34$ книги, на второй $52$ книги.
\end{enumerate}

\subsection*{К блоку 8}
\begin{enumerate}
  \item $4{,}7$ см.
  \item Острый; прямой; тупой; развёрнутый.
  \item $26$ см.
  \item $84$ см$^2$.
  \item Сторона $9$ см, периметр $36$ см.
  \item $74$ см$^2$.
  \item $9$ см.
  \item $72$ см$^3$.
  \item 1) $30\,000$ см$^2$; 2) $2500$ см$^3$.
  \item Ширина $11$ см, площадь $176$ см$^2$.
\end{enumerate}

\subsection*{К итоговой контрольной}
\begin{enumerate}
  \item $608\,000$.
  \item $672$.
  \item $210$.
  \item $\frac34$.
  \item $\frac34$.
  \item $11{,}75$.
  \item $x=60$.
  \item $160$ км.
  \item Длина $12$ см, периметр $40$ см.
  \item $49{,}5$ л.
\end{enumerate}

\section*{Финальная памятка перед началом 6 класса}
\addcontentsline{toc}{section}{Финальная памятка перед началом 6 класса}

\begin{center}
\begin{tabularx}{\textwidth}{@{}>{\bfseries}l X@{}}
\toprule
Навык & Ключевая идея \\
\midrule
Округление & смотрим на одну цифру справа от нужного разряда \\
Порядок действий & скобки $\to$ степени $\to$ умножение/деление $\to$ сложение/вычитание \\
Делимость & $2,5,10$ --- последняя цифра; $3,9$ --- сумма цифр \\
Сокращение дроби & делим числитель и знаменатель на общий делитель \\
Сложение дробей & сначала общий знаменатель \\
Деление дробей & умножаем на обратную дробь \\
Часть от целого & целое умножить на дробь \\
Целое по части & часть разделить на дробь \\
Десятичные дроби & контролируем разряды и положение запятой \\
Движение & $S=vt$ \\
Прямоугольник & $P=2(a+b)$, $S=ab$ \\
Параллелепипед & $V=abc$ \\
\bottomrule
\end{tabularx}
\end{center}

\vspace{1em}
\textcolor{accent}{\rule{\textwidth}{0.5pt}}

\textbf{Главный критерий готовности:} уметь выбрать наиболее подходящее правило в задаче, выполнить вычисления, проверить разумность результата и правильно записать единицы измерения.

\end{document}